
% Conceptos a explorar: Codominancia, Dominancia incompleta, Alelos múltiples, Pleiotropía, Alelos letales, Herencia ligada al sexo, Herencia poligénica

Si bien el trabajo de Mendel ayudó a explicar las bases de buena parte de la genética, un estudio limitado a una sola especie inevitablemente dejaría ciertos fenómenos fuera, esto dada la enorme diversidad de organismos en la naturaleza. Aun con esto, muchos de estos fenómenos pueden entenderse como variaciones de las bases mendelianas.

\section{Dominancia incompleta y codominancia}

Las primeras desviaciones de las leyes de Mendel surgen cuando la relación entre alelos no es de dominancia absoluta.

En la dominancia incompleta, el fenotipo del heterocigoto es una mezcla o un intermedio entre los fenotipos de los dos homocigotos. No existe un alelo dominante que oculte completamente al otro; ambos contribuyen al fenotipo, pero ninguno se expresa en su totalidad.

Definición: La dominancia incompleta es un tipo de herencia donde el fenotipo del heterocigoto es una combinación (generalmente un punto intermedio) de los fenotipos homocigotos.

Un ejemplo clásico es la planta de la boca de dragón (Antirrhinum majus). Al cruzar una planta de flores rojas (RR) con una de flores blancas (R'R'), todas las plantas de la generación F1 tienen flores rosas (RR'). Los alelos para el color rojo y blanco no son dominantes uno sobre el otro; en cambio, su interacción produce un fenotipo intermedio.

Codominancia

En la codominancia, ambos alelos de un gen se expresan completamente en el heterocigoto. En lugar de un fenotipo intermedio o mezcla, el individuo muestra ambos fenotipos de manera simultánea y distinguible.

Definición: La codominancia es un tipo de herencia donde ambos alelos se expresan por igual en el fenotipo del heterocigoto, sin que uno enmascare al otro.

Un ejemplo bien conocido en humanos es el sistema de grupos sanguíneos ABO. Los alelos IA e IB son codominantes. Un individuo con genotipo IAIB presenta ambos antígenos (A y B) en la superficie de sus glóbulos rojos, resultando en el tipo de sangre AB. Ni el antígeno A enmascara al B ni viceversa; ambos se expresan por igual.

Dato Biotecnológico: La codominancia es una herramienta fundamental en la identificación forense y en las pruebas de paternidad. Los marcadores genéticos codominantes, como los STRs (Repeticiones Cortas en Tándem), permiten distinguir entre todos los genotipos posibles (AA, AB, BB), lo que proporciona una alta capacidad de discriminación entre individuos.

\section{Alelos múltiples}

Mendel trabajó con genes que tenían dos variantes alélicas. Sin embargo, para muchos genes existen más de dos alelos posibles en una población. Esto se conoce como alelos múltiples.

Definición: Los alelos múltiples se refieren a la existencia de más de dos formas alélicas de un gen en una población. Un individuo diploide solo puede tener dos de estos alelos, pero la población en su conjunto puede presentar muchas variantes.

El sistema de grupos sanguíneos ABO en humanos es, de nuevo, un ejemplo excelente. Este gen tiene tres alelos principales: IA, IB e i. Las posibles combinaciones de estos tres alelos dan lugar a cuatro fenotipos: A (IAIA o IAi), B (IBIB o IBi), AB (IAIB) y O (ii). La presencia de alelos múltiples incrementa la diversidad genética de una población y es crucial para la adaptación.

\section{Pleiotropía}

Hasta ahora hemos considerado genes que afectan un solo rasgo. La pleiotropía describe el fenómeno opuesto: un único gen que influye en múltiples fenotipos aparentemente no relacionados.

Definición: La pleiotropía es la propiedad de un gen de afectar a múltiples características fenotípicas de un organismo.

Un ejemplo humano es la fenilcetonuria (PKU). Es una enfermedad metabólica hereditaria causada por una mutación en el gen que codifica la enzima fenilalanina hidroxilasa. Esta enzima es responsable de convertir el aminoácido fenilalanina en tirosina. Cuando la enzima es defectuosa, la fenilalanina se acumula en el organismo, causando una gama de síntomas que incluyen retraso mental, problemas de pigmentación (pelo y piel más claros) y un olor corporal característico. Un solo gen defectuoso produce múltiples efectos, demostrando la pleiotropía.

\section{Alelos letales}

Algunas mutaciones pueden ser tan perjudiciales que causan la muerte del organismo. Estos son los alelos letales. Pueden ser dominantes o recesivos.

Definición: Un alelo letal es una variante de un gen que provoca la muerte del organismo que lo porta. Pueden ser recesivos, causando la muerte solo en homocigosis, o dominantes, causando la muerte incluso en heterocigosis.

Un ejemplo histórico y clásico es el de los ratones de color amarillo. Un alelo dominante para el color del pelaje (Ay) produce un pelaje amarillo cuando está en heterocigosis (AyA). Sin embargo, cuando un ratón hereda dos copias de este alelo (AyAy), el embrión muere en el útero. Por lo tanto, todos los ratones amarillos vivos son heterocigotos, y el cruce de dos ratones amarillos produce una proporción fenotípica de 2 amarillos por 1 de otro color (en lugar de la esperada 3:1), ya que los homocigotos letales no llegan a nacer.

\section{Herencia ligada al sexo}

La herencia ligada al sexo se refiere a las enfermedades cuyos genes se localizan en los cromosomas sexuales X y Y, los cuales definen el sexo genético en el ser humano, XX para el femenino y XY para el masculino; por lo tanto, el único cromosoma común en ambos géneros es el X.\citep{delcastilloruizGeneticaClinica2019}

Los genes que se encuentran en los cromosomas sexuales (X e Y) presentan patrones de herencia particulares. En los humanos, la mayoría de los genes ligados al sexo se encuentran en el cromosoma X, ya que el Y es más pequeño y contiene menos genes.

Definición: La herencia ligada al sexo se refiere a la transmisión de genes localizados en los cromosomas sexuales. Los patrones de herencia para estos genes difieren entre machos y hembras.

Un ejemplo conocido es el daltonismo rojo-verde, una condición donde el individuo tiene dificultad para distinguir ciertos colores. El gen responsable está en el cromosoma X. Como los varones tienen un solo cromosoma X (heredado de la madre), si heredan un alelo mutante en ese X, serán daltónicos (herencia recesiva ligada al X). Las mujeres, al tener dos cromosomas X, necesitarían heredar el alelo mutante en ambos para presentar la condición, lo que es menos frecuente; si solo tienen un alelo mutante, son portadoras sanas.

\section{Herencia poligénica}

Los ejemplos anteriores suelen mostrar una variación fenotípica cualitativa y discreta (flor rosa o roja, sangre tipo A o B). Sin embargo, muchos de los rasgos más importantes son cuantitativos y muestran una variación continua (como la altura, el peso o el color de la piel). Esta variación continua es el resultado de la herencia poligénica.

Definición: La herencia poligénica es el tipo de herencia en la que un rasgo fenotípico está controlado por la acción combinada de múltiples genes, cada uno con un efecto pequeño y aditivo.

La altura humana es un ejemplo clásico. No es controlada por un solo gen, sino por la interacción de cientos de genes. Cada gen contribuye con una pequeña cantidad a la altura final. Además, el entorno (como la nutrición) también juega un papel crucial. El resultado es una distribución en forma de campana (normal) de la altura en la población, con la mayoría de los individuos alrededor de un promedio y menos individuos en los extremos (muy bajos o muy altos).

Dato Biotecnológico: Comprender la herencia poligénica es fundamental en la agricultura y la ganadería modernas. Los programas de mejora genética para aumentar la producción de leche en vacas, el rendimiento de grano en el maíz o la resistencia a enfermedades en plantas se basan en la selección de individuos con alelos favorables en múltiples genes poligénicos. La selección genómica, que utiliza marcadores moleculares en todo el genoma, es una herramienta biotecnológica clave para acelerar este proceso.